\section{Problem Setup}
\label{sec:setup}

\subsection{Task-agnostic inference over a shared vocabulary}

A frozen backbone with adapter parameters $\theta$ maps an input $\bx$ to a
hidden state $\bh_\theta(\bx)\in\R^d$. Decoding uses a frozen unembedding
$\bU\in\R^{V\times d}$ over the shared vocabulary. A classification-style task
$g$ is specified by a \emph{verbalizer} set $Y_g\subset[V]$ --- the tokens that
name its labels --- and its prediction is the restricted argmax
\begin{equation}
  \hat y_g(\bx)\;=\;\arg\max_{y\in Y_g}\; \bu_y^\top \bh_\theta(\bx).
  \label{eq:raw-decode}
\end{equation}
Tasks $g=1,\dots,T$ arrive in a stream; $\theta_g^{\mathrm{post}}$ denotes the
parameters just after task $g$ finishes training and $\theta_t$ the parameters at
a later time $t>g$. Nothing in \eqref{eq:raw-decode} is conditioned on $g$: at
inference the system sees $\bx$ and the option list, not the task index.

This is generative continual learning with a single decoder head, and it differs
from the vision continual-learning setting in one structural way that drives
everything below. There, each task owns a private classifier head, so re-fitting
that head is a per-task operation. Here there is one output space, and tasks
\emph{share tokens} inside it.

\subsection{Two capacity constraints, not one}

The literature budgets one resource: trainable parameters. A rank-$8$ LoRA bank
on T5-large $q/v$ costs $2.36$M scalars, and every method we compare against
spends its budget there. We separate a second resource.

For task $g$ at time $t$ define three retentions, all under the restricted argmax
of \eqref{eq:raw-decode}:
\begin{align}
  R_g^{\mathrm{raw}}(t) &= \mathrm{acc}_g\big(\theta_t,\ \text{no offset}\big),\\
  R_g^{\mathrm{orc}}(t) &= \max_{\bb_g\in\R^{|Y_g|}}
      \mathrm{acc}_g\big(\theta_t,\ \bb_g\big), \label{eq:orc}\\
  R_g^{\mathrm{shr}}(t) &= \mathrm{acc}_g\big(\theta_t,\ \bb^\star\big),
      \qquad \bb^\star \text{ one vector shared by all tasks}. \label{eq:shr}
\end{align}
An \emph{offset} is a vector $\bb\in\R^V$ added to the logits before the argmax,
$\hat y_g^{\,\bb}(\bx)=\arg\max_{y\in Y_g}(\bu_y^\top\bh_\theta(\bx)+\bb_y)$.
Equation~\eqref{eq:orc} allows a separate offset per task and therefore
\emph{requires task identity}; it is the generative analogue of the vision-CL
``re-fit the head'' quantity. Equation~\eqref{eq:shr} allows one offset for the
whole stream and requires nothing. By construction
$R_g^{\mathrm{raw}}\le R_g^{\mathrm{shr}}\le R_g^{\mathrm{orc}}$.

Work on shallow versus deep forgetting studies the outer gap
$R_g^{\mathrm{orc}}-R_g^{\mathrm{raw}}$ and, finding it large, declares it cheap:
a linear probe recovers it. We study the inner gap
\begin{equation}
  \Delta_g^{\mathrm{id}}(t)\;=\;R_g^{\mathrm{orc}}(t)-R_g^{\mathrm{shr}}(t)\;\ge\;0,
  \label{eq:delta-id}
\end{equation}
the \emph{identification gap}: the part of shallow forgetting that is not
recoverable without knowing which task a query came from. In the vision-CL
framing $\Delta^{\mathrm{id}}$ is zero by assumption, because a private head per
task \emph{is} a per-task offset. Under task-agnostic generative inference it is
not zero, and it is not addressed by spending more parameters.

\paragraph{Why parameter count is the wrong budget for this term.}
An offset supported on the union of all verbalizers costs
$|\bigcup_g Y_g|$ scalars --- at most $210$ on the $15$-task stream of
Section~\ref{sec:experiments}, against $2.36$M for the adapter bank, i.e.\
under $0.01\%$. The offset is parametrically negligible. The content of this
paper is that being parametrically negligible does not make it free: the binding
constraint on $\Delta^{\mathrm{id}}$ is not parameter count but
\emph{identifiability} --- how many distinct offsets a task-free inference rule
can select between, and how well those selections match what each task wanted.

\subsection{Scopes: the structure that makes the gap nonzero}

Call two tasks \emph{scope-mates} when they present the same label option set,
and let the \emph{scope} $S$ of a task be that set. A scope is exactly what an
inference-time rule can read off a prompt without a task index. Write $K_S$ for
the number of labels in $S$ and $\mathrm{mem}(S)$ for its member tasks.

\begin{proposition}[Disjoint verbalizers give a zero gap]
\label{prop:disjoint}
If the sets $Y_g$ are pairwise disjoint then a single shared $\bb$ realises every
per-task optimum simultaneously, so $\Delta_g^{\mathrm{id}}(t)=0$ for all $g,t$.
\end{proposition}

\noindent We state Proposition~\ref{prop:disjoint} first because it is the honest
boundary of the whole phenomenon: the identification gap requires tasks that
\emph{reuse output tokens}. It is also a falsifiable wiring check --- on any
singleton scope the measured $\Delta^{\mathrm{id}}$ must be exactly zero --- and we
report it as such in every experiment below.

The precondition is a property of the benchmark, not an assumption. On the
$15$-task stream we use, the shared-verbalizer groups are
\begin{center}
\small
\begin{tabular}{llc}
\toprule
label set & member tasks & $K_S$ \\
\midrule
\texttt{\{bad, good\}}                 & IMDB, SST-2 & 2 \\
\texttt{\{false, true\}}               & WiC, QQP, BoolQA, MultiRC & 2 \\
\texttt{\{contradiction, entailment, neutral\}} & MNLI, CB & 3 \\
\texttt{\{very negative, \ldots, very positive\}} & Yelp, Amazon & 5 \\
\bottomrule
\end{tabular}
\end{center}
plus five singleton scopes. Ten of the fifteen tasks share output tokens with at
least one other task, so $\Delta^{\mathrm{id}}$ is a live quantity on this stream
rather than a hypothetical one.

\subsection{The identification ladder}
\label{sec:ladder}

Between the extremes of \eqref{eq:orc} and \eqref{eq:shr} sits a ladder a
task-free system can actually climb. Let $\bb_g^\star$ be an optimal per-task
offset, gauge-fixed to zero mean on its own scope --- adding a constant to every
logit of a scope leaves the restricted argmax unchanged, so only the zero-mean
part of an offset is identified. Four rungs, in increasing order of what they
require at inference time:
\begin{enumerate}
  \item $R^{\mathrm{raw}}$: no offset.
  \item $R^{\mathrm{shr}}$: one offset for the entire stream. Requires nothing.
  \item $R^{\mathrm{shr}}_{\mathrm{scope}}$: one offset per scope, selected by
        reading the option list. Requires no task index.
  \item $R^{\mathrm{orc}}$: one offset per task. Requires the task index and is
        therefore an upper bound, never a deployable method.
\end{enumerate}
Rung 3 is the interesting one, and Proposition~\ref{prop:disjoint} tells us why:
on a singleton scope it coincides with rung 4, so any gap it leaves lives
entirely on multi-task scopes, where scope-mates disagree about which direction
their verbalizer logits should move. A finer rung interpolates 3 and 4: allow
$m$ offsets per scope and select among them by a task-free rule, giving
$R^{\mathrm{book}}(m)$ with $R^{\mathrm{shr}}_{\mathrm{scope}}=R^{\mathrm{book}}(1)$.
The quantization radius of that $m$-centre codebook, $q_m(S)$, is the object our
theory bounds and our experiments measure.

\paragraph{Budget accounting we hold ourselves to.}
Three resources, reported for every arm, because a comparison is capacity-matched
only if all three match: (i) trainable parameters deployed; (ii) stored state at
inference, counted in scalars, including anything an offset table holds; and
(iii) whether any old-task example is re-read at any point. Item (iii) is not a
footnote. A per-scope offset refit at each stage from all seen tasks' data is
rehearsal, and comparing it against a rehearsal-free method as though both were
deployable is the error this paper's experimental section is organised to avoid.
